\chapter{Diskussion und Bewertung}
\label{ch:diskussion}

Dieses Kapitel bewertet die entwickelte Systemarchitektur hinsichtlich ihrer technischen Eignung, Robustheit und Übertragbarkeit. Dabei werden sowohl algorithmische als auch systemtechnische Aspekte analysiert.

\section{Bewertung der gewählten Gesamtarchitektur}

Die Entscheidung für eine strikt scan-basierte Pipeline stellt eine bewusste Abkehr von komplexeren, mehrstufigen Wahrnehmungsansätzen dar. Diese Reduktion des Datenpfads führt zu mehreren signifikanten Vorteilen.

\subsection{Reduktion von Komplexität}

Durch den direkten Übergang von \texttt{/scan} zu einer gefilterten Scan-Repräsentation wird auf zusätzliche Abstraktionsebenen wie Objektextraktion oder Hindernismodellierung verzichtet. Dies reduziert:

\begin{itemize}
    \item die Anzahl potenzieller Fehlerquellen,
    \item den Parametrierungsaufwand,
    \item die Latenz innerhalb der Verarbeitungskette.
\end{itemize}

Gerade für reaktive Verfahren wie Follow-the-Gap ist diese Vereinfachung fachlich sinnvoll, da der Algorithmus explizit für rohnahe Sensordaten formuliert wurde~\cite{sezer2012}.

\subsection{Vergleich mit alternativen Ansätzen}

Im Gegensatz zu globalen Planern oder kartengestützten Verfahren besitzt der implementierte Stack keine Kenntnis über die Umgebung außerhalb des aktuellen Sichtfelds. Dies führt zu folgenden Konsequenzen:

\begin{itemize}
    \item hohe Reaktionsfähigkeit auf lokale Hindernisse,
    \item jedoch eingeschränkte Fähigkeit zur global optimalen Pfadplanung.
\end{itemize}

Die Architektur ist daher klar als \textbf{lokaler reaktiver Planer} einzuordnen und nicht als vollständiges Navigationssystem.



\section{Analyse der Planner- und Control-Schicht}

Die Trennung zwischen algorithmischem Kern, Planner und Control stellt eine klare funktionale Hierarchie dar.

Die gewählte Aufteilung ermöglicht:

\begin{itemize}
    \item Wiederverwendbarkeit des FTG-Kerns,
    \item unabhängige Anpassung der Fahrstrategie,
    \item klare Definition von Verantwortlichkeiten innerhalb der Pipeline.
\end{itemize}

Insbesondere fungiert der Planner als \textbf{Policy-Schicht}, die sicherstellt, dass die reaktive Planung in physikalisch realisierbare und sichere Fahrbefehle überführt wird.


\section{Bewertung der experimentellen Ergebnisse}

Die durchgeführten Fahrversuche zeigen, dass das System in strukturierten Innenräumen grundsätzlich funktionsfähig ist, jedoch klare Abhängigkeiten von der Umgebungskonfiguration aufweist.

\subsection{Quantitativer Vergleich der Validierungsläufe}

Tabelle~\ref{tab:validierung_vergleich} fasst die wichtigsten Metriken der beiden Testläufe zusammen.

\begin{table}[H]
\centering
\caption{Vergleich der Metriken: Validierung Test 1 vs. Test 2}
\label{tab:validierung_vergleich}
\begin{tabular}{l c c}
\toprule
\textbf{Metrik} & \textbf{Test 1} & \textbf{Test 2} \\
\midrule
Gap erkannt & 90\% & 56\% \\
Speed = 0 (Zeitanteil) & 10\% & 44\% \\
$\emptyset$ Geschwindigkeit & 0,37 m/s & 0,40 m/s \\
Min. Clearance & 0,27 m & 0,34 m \\
Kollisionen & 0 & 0 \\
Heading-Bereich & $\pm$ 0,8 rad & $\pm$ 0,8 rad \\
\midrule
\bottomrule
\end{tabular}
\end{table}

\subsection{Detailbewertung der Szenarien}

\begin{itemize}
    \item \textbf{Test 1:} Alle Hindernisse wurden umfahren und die 90°-Ecke gemeistert. Der kurze Stopp an der Ecke resultiert aus dem \texttt{kCarRadius}-Parameter von 0,4\,m, der hier eine konservative Reaktion auslöste.
    \item \textbf{Test 2:} Dieser Lauf ist als eingeschränkt zu bewerten. Ein Zeitanteil von 44\% mit Geschwindigkeit 0 sowie eine Gap-Erkennungsrate von nur 56\% weisen auf Instabilitäten hin. Die Ursache liegt im 180° FOV, das zu breit gewählt war, wodurch Seitenwände die Lückenerkennung störten.
\end{itemize}

\subsection{Empfehlungen aus der Validierung}
Bewertung für alle Parameter.
Basierend auf den Ergebnissen wird empfohlen, das FOV auf 120° zurückzusetzen, um Störeinflüsse durch Wände zu minimieren. Zudem sollte der \texttt{kCarRadius} von 0,4\,m auf 0,20\,m reduziert werden, um dem realen 30\,cm breiten Fahrzeug besser gerecht zu werden.

\subsection{Robustheit gegenüber statischen Hindernissen}

In beiden Versuchsreihen konnte das Fahrzeug Hindernisse erkennen und geeignete Ausweichbewegungen berechnen. Die Anpassung der Geschwindigkeit in Abhängigkeit vom Frontabstand erwies sich als stabil und reproduzierbar.

\subsection{Verhalten bei dynamischen Änderungen}

Besonders relevant ist das Verhalten bei der Entfernung eines Hindernisses während der Fahrt. Hier zeigt sich ein charakteristischer Vorteil reaktiver Systeme:

\begin{itemize}
    \item sofortige Anpassung an die neue Umgebung,
    \item keine Notwendigkeit zur Replanung einer globalen Route.
\end{itemize}

Dieses Verhalten bestätigt die Eignung des FTG-Ansatzes für dynamische Szenarien.

\subsection{Grenzfälle}

Trotz der positiven Ergebnisse treten typische Einschränkungen auf:

\begin{itemize}
    \item keine strategische Entscheidung in Sackgassen,
    \item potenziell suboptimale Pfadwahl in komplexen Geometrien,
    \item Abhängigkeit von Sichtfeldbegrenzungen.
\end{itemize}

Diese Einschränkungen sind inhärent im Algorithmus begründet und nicht als Implementierungsfehler zu interpretieren.

\section{Einordnung in den wissenschaftlichen Kontext}

Die Kombination aus klassischem FTG-Ansatz~\cite{sezer2012}, CTU-Referenzimplementierung~\cite{ctu_ftg} und moderner ROS~2-Infrastruktur~\cite{ros2_docs} zeigt, wie etablierte Verfahren in aktuelle Softwarearchitekturen überführt werden können.

Die vorliegende Arbeit leistet dabei keinen Beitrag zur Weiterentwicklung des Algorithmus selbst, sondern zur:

\begin{itemize}
    \item systematischen Integration,
    \item reproduzierbaren Umsetzung,
    \item technischen Konsolidierung eines bestehenden Ansatzes.
\end{itemize}

